\section{Quantum evolution along a genuine growing trace}
\label{sec:growing-trace}

The first task in the physical program is to determine the quantum
evolution induced by a specified Loewner curve within a fixed theory.
This does not require an explicit expectation value or an immediate
identification with the arithmetic transfer. We now carry out that
calculation for an independently specified observable in pure
four-dimensional Yang--Mills theory. Unlike the regular inverse-map
deformations in Section~\ref{sec:realization}, the path here is the
actual growing trace; previously traced points remain fixed.

\subsection{A fixed theory and a specified tip contraction}

Fix Euclidean $SU(N)$ Yang--Mills at $\theta=0$, including its
coupling and state. Embed the upper half-plane in $x_3=x_4=0$ and
let a deterministic driver $u(t)$ generate a sufficiently regular
trace $q[0,t]$, starting at $q(0)=0$. The slit is geometric data,
not a moving physical boundary of the quantum theory. The driver
need not be random for the field expectation to be quantum.

An open holonomy carries color indices at its endpoints. We
contract them by returning from $q(t)$ to the origin along the
straight chord. This return path is part of the observable. Let
$U_t$ transport along the actual prefix and $R_t(r)$ transport
along the straight segment from $0$ to $r q(t)$. Define
\begin{equation}\label{eq:main-trace-observable}
 \mathsf H_t=R_t(1)^{-1}U_t,\qquad
 W_\tau(t)=\frac1N\vev{\tr\mathsf H_t},\qquad W_\tau(0)=1.
\end{equation}
To state smooth-field identities with a controlled short-time
expansion, the connection is the Yang--Mills flowed field $B_\tau$
at a fixed positive resolution $\tau$. The action and measure
remain those of the original interacting theory; there is no
Gaussian substitution. Field-flow time $\tau$ is distinct from
Loewner time $t$ and the arithmetic shift $\omega$
\cite{LuscherFlow,LuscherWeiszFlow}.

The following identities are exact for smooth connections. Their
expectation-value versions assume existence of the selected
regulated or flowed correlators and the integrability needed to
differentiate them. No nonperturbative continuum construction is
asserted. Removing the resolution requires a separate treatment
of cusp and contact terms. Neither this observable nor its
normalization is identified with $V_{\omega,L}$.

\subsection{The tip cancellation and the quantum hierarchy}

Suppress $\tau$ in field notation and transport the curvature
back to the fixed basepoint:
\begin{equation}\label{eq:main-trace-moment}
 \widehat F_t(r)=R_t(r)^{-1}F_{12}(r q(t))R_t(r),\quad
 \mathsf J_t=\int_0^1r\widehat F_t(r)\dd r,\quad
 k_t=q_1\dot q_2-q_2\dot q_1.
\end{equation}
Prefix growth inserts the connection at the tip. Differentiating
the return chord produces the opposite endpoint term, so these
terms cancel before averaging. The remaining exact matrix and
expectation equations are
\begin{equation}\label{eq:main-trace-flow}
 \boxed{\dot{\mathsf H}_t=\ii k_t\mathsf J_t\mathsf H_t,
 \qquad \dot W_\tau(t)=\ii k_t M_1(t),\qquad
 M_1=\frac1N\vev{\tr(\mathsf J_t\mathsf H_t)}.}
\end{equation}
The geometric coefficient is twice the oriented area growth of
the trace with its chord. The quantum coefficient is a curvature
insertion in the same interacting expectation.

Differentiating once more gives
\begin{equation}\label{eq:main-trace-second}
 \dot M_1=\dot q^\mu M_{D,\mu}+2\ii k_t M_2,
 \qquad \mu=1,2.
\end{equation}
Here $M_{D,\mu}$ is the normalized expectation with insertion
$\int_0^1r^2\widehat{D_\mu F_{12}}_t(r)\dd r$; $M_2$ has the
ordered insertion
\[
 \int_{0<s<r<1}rs\widehat F_t(s)\widehat F_t(r)\dd s\dd r.
\]
Both are followed by $\mathsf H_t$ inside the trace. The order and
the factor of two follow from the changing radial transports.
The complete derivation and domains are given in
\cite[Section~S13]{Supplement}. These are equations for quantum
expectations without a Wick rule or large-$N$ factorization;
they are not yet a closed scalar equation for $W_\tau$.

\subsection{A linear driver and a controlled short-time response}

For the concrete driver $u(t)=a t$, the upper-half-plane tip obeys
\begin{equation}\label{eq:main-linear-trace}
 aq+2\log(1-aq/2)=a^2t,\qquad
 q(t)=2\ii\sqrt t+\frac{2a}{3}t-\frac{\ii a^2}{18}t^{3/2}
                       +O_a(t^2).
\end{equation}
The implicit relation is for $a\ne0$ on the branch issuing from
the origin; its $a=0$ limit is the vertical slit. It follows by
integrating the forward Loewner equation, not by evaluating that
equation at its singular tip. The enclosed oriented area is
$\mathcal A(t)=-(2a/9)t^{3/2}+O_a(t^{5/2})$. Under the local
moment assumptions stated in the companion,
\begin{equation}\label{eq:main-trace-short-time}
 \boxed{W_\tau(t)=1-\frac{2a^2}{81}C_\tau t^3
                         +O_{\tau,a}(t^{7/2}),\qquad
 C_\tau=\frac1N\vev{\tr F_{12}[B_\tau](0)^2}.}
\end{equation}
The coefficient $C_\tau$ is a gauge-invariant moment in the fixed
theory. Its value need not be known to derive this relation.
The companion bounds the remainder explicitly by curvature and
covariant-derivative moments; the expansion is at fixed $\tau$.
For the constant driver the chosen loop backtracks exactly,
so $\mathsf H_t=I$ and $W_\tau(t)=1$ configuration by configuration.
That control concerns this return-path choice, not every Wilson
observable on a straight trace.

\subsection{The Makeenko--Migdal closure question}

Schwinger--Dyson identities relate field-equation insertions to
derivatives of observables. Their loop-equation form is a natural
candidate for reducing \eqref{eq:main-trace-second}
\cite{MakeenkoLoops}. The first missing terms are explicit. In
four dimensions $D_1F_{12}$ is only one part of
$E_2=\sum_\mu D_\mu F_{\mu2}$: the terms $D_3F_{32}$ and
$D_4F_{42}$ remain even though the contour is planar. Finite-$N$
splitting terms also retain expectations of products. At positive
$\tau$, differentiating the observable with respect to the
original field differentiates the flow map $B_\tau[A]$ as well.
These terms have not been reduced to $W_\tau$ and $M_1$.

Thus the calculation supplies a concrete relation between genuine
Loewner growth and quantum evolution, and identifies the
correlators needed for a closure attempt. It supplies neither a
closed Makeenko--Migdal reduction nor the physical arithmetic
transfer. The freedom to choose another return path, endpoint
contraction or insertion stays within the fixed-theory question;
each choice must have its own derived evolution.
